% FRAGMENT A06 (drafter A: class model and main results)
% REPLACES: main.tex lines 264–384 ("\subsection{2.2 Values, preparation, and the" to "… \(B_r(\mu)=g_L+\mu(g_H-g_L)\).")
%   This covers Sections 2.2, 2.3, and 2.4.
% ANCHORS FOR OTHER FRAGMENTS (LaTeX comments in the body):
%   - "ANCHOR M1_03 BEGIN/END": M1_03 replaces the text between these two comments
%     (it is the original text of main.tex lines 337–340, kept here verbatim).
%   - "ANCHOR M1_02": insert M1_02 (roles paragraph) at this comment, that is, after the
%     sentence that ends "… depends on the price through preparation." (old line 349).
% LABELS KEPT: sec-values-preparation-and-the-sale, eq:paper-1, sec-trading-and-timing,
%   eq:paper-2, eq:paper-3, sec-equilibrium.
% LABELS REMOVED: none.  LABELS ADDED: none.
% NEW SYMBOLS: j (quality class), \pi (prior), G_j (class value cdf), G_j(p^-),
%   D(u)=G_L(u)-G_H(u), \bar r (upper end of the incumbent support), t(\theta), g(\theta),
%   B_F, g_j(F), \Delta_T(F) (objects under incumbent law F).
% REMOVED SYMBOLS: c_L (the low cost is zero), B_r as the general symbol (B_r is now the
%   uniform special case of B_F).
% NEW CITATIONS: none.  NUM PLACEHOLDERS: none.
% EQUATION TAGS: (1), (2), (3) unchanged. New displays are unnumbered.

\subsection{2.2 Values, preparation, and the
sale}\label{sec-values-preparation-and-the-sale}

The target's known standalone value is normalized to zero. Acquisition
values and preparation costs are measured per target share. Adding the
same standalone component to every ownership outcome shifts prices and
payoffs by a constant without changing incentives.

All strategic agents are risk neutral. In the benchmark the target
share is the only traded claim with information about the acquisition
match, and neither bidder trades it. The investor has no control
rights and cannot acquire the target. The sale mechanism binds all
target shares. Shareholder tendering and holdout stay outside the
model. The incumbent's value distribution is public, conditional on
the information at the time the sale opportunity becomes visible. Its
realized value stays undisclosed to the challenger before preparation.

Two potential acquirers face the target. The incumbent has already
prepared and incurs no further participation cost. Its acquisition value
\(R\) has a continuous distribution \(F\) on \([0,\bar r]\), conditional
on public information when the sale process begins, and \(F\) is
extended by unity above \(\bar r\). A first-order stochastic shift of
\(F\) strengthens the incumbent. In the leading case \(R\) is uniform on
\([0,r]\), and an increase in \(r\) strengthens the incumbent. The
incumbent is a bidder, rather than the target's management, and knows
its value before bidding.

The challenger belongs to one of two quality classes, \(j\in\{H,L\}\),
with prior \(\pi=\Pr(j=H)\in(0,1)\). Given its class, its acquisition
value \(\theta\) has a distribution \(G_j\) on \([0,\infty)\) with a
finite mean; atoms are allowed. Write \(G_j(p^-)\) for the probability
that a class-\(j\) value falls below the reserve \(p\), and
\(D(u)=G_L(u)-G_H(u)\). Two restrictions rank the classes. The high
class dominates the low class in the sense of first-order stochastic
dominance above the reserve, and the two classes are equally likely to
fall below the reserve, which I call equal admissibility:

\[
G_H(u)\le G_L(u)\quad\text{for every }u>p,\qquad
G_H(p^-)=G_L(p^-).
\]

The simplest case has both supports above \(p\). With no atom at the
reserve in either class, the two restrictions are exactly the condition
for competition to widen the class difference in target proceeds under
every incumbent distribution
(\hyperref[result-proposition-1]{Proposition 1}). The structure is the
selective-entry structure of the auction literature. The class indexes
the distribution of the challenger's value, preparation reveals the
value, and the stock price is the pre-entry signal.

The benchmark is the point-mass case: \(G_H\) puts all its mass on
\(h\), \(G_L\) puts all its mass on \(\ell\), and \(\pi=1/2\). In the
benchmark the class and the value coincide, and values, the reserve, and
the uniform incumbent's strength satisfy

\[
0<p<\ell<r<h.
\tag{1}\label{eq:paper-1}
\]

The general statements need no normalization of money. A common
positive rescale of every value, the reserve, and every preparation
and trading cost rescales prices and payoffs by the same factor.
Orders, beliefs, entry, and ownership do not change. The benchmark
measures money in units of the low value, so \(\ell=1\) there.

The challenger initially knows neither its class nor its value. After
observing the stock price, it privately learns its preparation cost
\(C\in\{0,c_H\}\), with \(\Pr(C=0)=\rho\), \(0<\rho<1\), and
\(c_H>0\). The other type's cost is zero: a low-cost
challenger is a ready buyer whose diligence costs nothing at the margin.
There is one challenger with one realized cost, not a low-cost buyer and
a high-cost buyer at the same time; \(\rho\) is the probability that the
realized cost is zero. Paying \(C\) reveals \(\theta\) and permits
bidding; declining leaves the challenger outside the sale. The class and
value of the challenger, the incumbent's value, the preparation cost,
and noise demand are mutually independent, so the cost carries no
information about the class or the value.

Before trading, the seller publicly commits to a cash second-price
auction with reserve \(p>0\). The highest admissible bidder acquires the
target and pays the larger of the reserve and the highest competing
admissible bid. There is no sale without an admissible bid. Both bidders
use truthful, weakly dominant bids, and a prepared challenger whose
value is below the reserve does not bid. A bid equal to the reserve is
admissible, and a challenger indifferent about preparation enters.
Because \(F\) is continuous, acquisition-value ties have probability
zero. Preparation indifference matters at a posterior plateau considered
in Section 4.

\subsection{2.3 Trading and timing}\label{sec-trading-and-timing}

An investor observes the challenger's class \(j\) and submits an order
\(q\in[-1,1]\). The bound is a normalized small trading unit, not
ownership of the target. The investor has no initial position and pays
a linear trading cost \(k|q|\), where \(k>0\). This cost is a separate
friction, such as an execution or position-carrying cost. It is
distinct from adverse-selection price impact, which the competitive
pricing rule below already generates. Its profit and aggregate order
flow are

\[
\begin{gathered}
q\{V_T-P(X)\}-k|q|,\qquad X=q+Z,\\
f(z)=\frac{1}{2b}e^{-|z|/b},\qquad b>1.
\end{gathered}
\tag{2}\label{eq:paper-2}
\]

% ANCHOR M1_03 BEGIN: fragment M1_03 replaces the text from here to "ANCHOR M1_03 END".
Here \(V_T\) is the terminal payoff of a target share and \(Z\) is
independent noise demand with a Laplace density. The bounded likelihood
ratios of this density limit the information that any order can reveal.
\hyperref[sec-extensions]{Section 5} also considers logistic noise.
% ANCHOR M1_03 END.

Competitive market makers observe aggregate flow and set

\[
P(X)=\mathbb E[V_T\mid X],
\tag{3}\label{eq:paper-3}
\]

anticipating preparation and auction outcomes. The trading stage follows the single-auction model of Kyle (\citeproc{ref-Kyle1985}{1985}), with binary information, bounded orders, a linear trading cost, Laplace noise, and a terminal payoff that depends on the price through preparation.

% ANCHOR M1_02: insert fragment M1_02 (roles paragraph) here, as a new paragraph.

The information sets are as follows. Market makers observe aggregate
flow \(X\) and nothing else. The challenger observes the price \(P\) and
its own cost \(C\), but not \(X\), its class, its value \(\theta\), or
the incumbent's realized value; in the extension of
\hyperref[sec-complementary-private-information]{Section 5.3} it also
observes a private signal, which the investor does not see. The investor
observes the class \(j\) but not the value \(\theta\), and nothing about
the incumbent's realized value beyond its public distribution. The
incumbent knows its own value. The sequence is: the seller announces the
sale rule; the investor learns the class and trades; market makers set
the price from flow; the challenger observes the price and its cost and
decides whether to prepare; prepared bidders learn their values and
submit truthful bids; ownership and financial payoffs are realized.

\subsection{2.4 Equilibrium}\label{sec-equilibrium}

An equilibrium specifies class-contingent order distributions
\(\sigma_H,\sigma_L\), a measurable price function, Bayesian beliefs
conditional on the observed price, optimal preparation, and truthful
bidding. The investor can mix and can deviate to any order in
\([-1,1]\). Two comparisons recur. A unilateral deviation faces a fixed
candidate price function and preparation schedule, and changes only
the distribution of flow that reaches them. A comparison across
economies, such as across incumbent strengths, re-solves the price
function and the preparation schedule in each economy. Proofs use the
first comparison to test a candidate equilibrium and the second to
compare outcomes.

For a known challenger value \(\theta\), write \(t(\theta)\) for
expected target proceeds after the challenger prepares and
\(g(\theta)\) for the challenger's expected gross acquisition profit;
\(t_0\) is expected target proceeds without entry. The class objects are
expectations over the class distribution,

\[
t_j=\mathbb E_j[t(\theta)],\qquad g_j=\mathbb E_j[g(\theta)],\qquad
j\in\{H,L\},
\]

and the target-payoff spread is \(\Delta_T=t_H-t_L\). Bids are truthful,
so these objects depend on the incumbent's distribution and the reserve
but not on any belief. At a belief \(\mu=\Pr(j=H)\), the challenger's
expected gross profit is \(B(\mu)=g_L+\mu(g_H-g_L)\), because its value
depends on the price only through its class. When the incumbent's
distribution matters I write \(g_j(F)\), \(\Delta_T(F)\), and \(B_F\);
for the uniform incumbent on \([0,r]\) I write \(g_j(r)\),
\(\Delta_T(r)\), and \(B_r\).
